\uuid{uI4M}
\titre{Introduction aux séries de Fourier}
\niveau{L2}
\module{Analyse}
\chapitre{Série de Fourier}
\sousChapitre{Calcul de coefficients}
\theme{Orthogonalité, approximation trigonométrique, série de Leibniz}
\auteur{Maxime Nguyen}
\datecreate{2026-05-19}
\organisation{AMSCC}
\difficulte{2}

\contenu{
	
	\texte{
		Peut-on reconstruire une fonction périodique à partir des fonctions
		$$
		1,\quad \cos(x),\quad \sin(x),\quad \cos(2x),\quad \sin(2x),\quad \ldots ?
		$$
		On considère la fonction $f$, $2\pi$-périodique, définie sur $[-\pi,\pi]$ par $f(x)=x$.
	}
	
	\begin{enumerate}
		
		%-------------------------------------------------------
		\item 
			

			On cherche à l'approcher par une fonction de la forme $S_1(x)=a\cos(x)+b\sin(x)$.}
		
		\begin{enumerate}
			
			\item \question{Représenter graphiquement la fonction $f$ sur $[-\pi,\pi]$.}
			\reponse{La fonction $f : x \mapsto x$ est la droite passant par l'origine, de pente $1$,
				tracée sur $[-\pi,\pi]$.}
			
			\item \question{La fonction $f$ est-elle paire ou impaire ?}
			\reponse{Pour tout $x$, $f(-x) = -x = -f(x)$, donc $f$ est impaire.}
			
			\item \question{Parmi $\cos(x)$ et $\sin(x)$, laquelle semble mieux adaptée pour approcher
				$f$ ? Que peut-on alors conjecturer sur le coefficient $a$ ?}
			\reponse{$\cos(x)$ est une fonction paire, tandis que $\sin(x)$ est impaire, comme $f$.
				Approcher une fonction impaire par une combinaison de fonctions implique que les termes pairs
				n'ont pas leur place : on conjecture $a = 0$.}
			
		\item \question{On choisit donc une approximation de la forme $S_1(x)=b\sin(x)$.
		On souhaite choisir $b$ de sorte que l'erreur quadratique intégrale
	\[
	E(b) = \int_{-\pi}^{\pi}\bigl(f(x)-b\sin(x)\bigr)^2\,dx
	\]
	soit minimale. Montrer que minimiser $E(b)$ revient à imposer la condition
	\[
	\int_{-\pi}^{\pi}\left(f(x)-b\sin(x)\right)\sin(x)\,dx = 0.
	\]}
\reponse{$E(b)$ est une fonction polynomiale de degré 2 en $b$, à coefficients réels.
	En développant :
	\[
	E(b)
	= \int_{-\pi}^{\pi} f(x)^2\,dx
	- 2b\int_{-\pi}^{\pi} f(x)\sin(x)\,dx
	+ b^2\int_{-\pi}^{\pi}\sin^2(x)\,dx.
	\]
	Comme $\int_{-\pi}^{\pi}\sin^2(x)\,dx = \pi > 0$, $E$ est une parabole ouverte vers le haut,
	et admet un unique minimum atteint en $b^*$ vérifiant $E'(b^*)=0$, soit
	\[
	-2\int_{-\pi}^{\pi}f(x)\sin(x)\,dx + 2b^*\int_{-\pi}^{\pi}\sin^2(x)\,dx = 0,
	\]
	ce qui est exactement la condition d'orthogonalité
	$\displaystyle\int_{-\pi}^{\pi}\bigl(f(x)-b^*\sin(x)\bigr)\sin(x)\,dx = 0$.}
		\end{enumerate}
		
		%-------------------------------------------------------
		\item \texte{\textbf{Comment choisir le « meilleur » coefficient ?}
			
			Pour déterminer $b$, on impose que l'erreur soit orthogonale à $\sin(x)$, c'est-à-dire :
			\[
			\int_{-\pi}^{\pi} \bigl(f(x)-b\sin(x)\bigr)\sin(x)\,dx=0.
			\]}
		
		\begin{enumerate}

			
			\item \question{Calculer $\displaystyle\int_{-\pi}^{\pi} \sin^2(x)\,dx$.}
			\reponse{En utilisant la formule $\sin^2(x)=\dfrac{1-\cos(2x)}{2}$ :
				\[
				\int_{-\pi}^{\pi} \sin^2(x)\,dx
				= \int_{-\pi}^{\pi} \frac{1-\cos(2x)}{2}\,dx
				= \frac{1}{2}\Bigl[x - \frac{\sin(2x)}{2}\Bigr]_{-\pi}^{\pi}
				= \pi.
				\]}
			
			\item \question{Calculer, à l'aide d'une intégration par parties,
				$\displaystyle\int_{-\pi}^{\pi} x\sin(x)\,dx$.}
			\indication{Poser $u=x$ et $v'=\sin(x)$.}
			\reponse{On pose $u=x$, $v'=\sin(x)$, donc $u'=1$, $v=-\cos(x)$.
				\begin{align*}
					\int_{-\pi}^{\pi} x\sin(x)\,dx
					&= \Bigl[-x\cos(x)\Bigr]_{-\pi}^{\pi} + \int_{-\pi}^{\pi}\cos(x)\,dx \\
					&= \bigl(-\pi\cos(\pi)+(-\pi)\cos(-\pi)\bigr) + \Bigl[\sin(x)\Bigr]_{-\pi}^{\pi} \\
					&= \bigl(\pi + (-\pi)(-1)\bigr) \cdot (-1)\cdot(-1) + 0\\
					&= \pi \cdot 1 + \pi \cdot 1 = 2\pi.
				\end{align*}
				Plus directement : $[-x\cos x]_{-\pi}^{\pi} = -\pi(-1)-(-\pi)(-1)(-1)$\ldots
				On obtient $\displaystyle\int_{-\pi}^{\pi} x\sin(x)\,dx = 2\pi$.}
			
			\item \question{En déduire la valeur de $b$.}
			\reponse{On a $b \cdot \pi = 2\pi$, donc $b = 2$.
				La première approximation est $S_1(x) = 2\sin(x)$.}
			
		\end{enumerate}
		
		%-------------------------------------------------------
		\item \texte{\textbf{Améliorer l'approximation.}
			
			On cherche maintenant une approximation plus précise de $f$, sous la forme
			\[
			S_2(x)=b_1\sin(x)+b_2\sin(2x).
			\]
			On impose que l'erreur $f(x)-S_2(x)$ soit orthogonale à $\sin(x)$ et à $\sin(2x)$.}
		
		\begin{enumerate}
			
			\item \question{Écrire les deux conditions intégrales correspondantes.}
			\reponse{
				\[
				\int_{-\pi}^{\pi} \bigl(f(x)-b_1\sin(x)-b_2\sin(2x)\bigr)\sin(x)\,dx = 0,
				\]
				\[
				\int_{-\pi}^{\pi} \bigl(f(x)-b_1\sin(x)-b_2\sin(2x)\bigr)\sin(2x)\,dx = 0.
				\]}
			
			\item \question{Montrer que
				$\displaystyle\int_{-\pi}^{\pi} \sin(x)\sin(2x)\,dx=0$.}
			\indication{Utiliser la formule $\sin(x)\sin(2x) = \tfrac{1}{2}[\cos(x)-\cos(3x)]$.}
			\reponse{Par la formule produit-en-somme :
				\[
				\sin(x)\sin(2x) = \frac{1}{2}\bigl[\cos(x-2x)-\cos(x+2x)\bigr]
				= \frac{1}{2}\bigl[\cos(x)-\cos(3x)\bigr].
				\]
				Donc
				\[
				\int_{-\pi}^{\pi}\sin(x)\sin(2x)\,dx
				= \frac{1}{2}\Bigl[\sin(x)-\frac{\sin(3x)}{3}\Bigr]_{-\pi}^{\pi} = 0.
				\]
				C'est un cas particulier du fait que $(\sin(nx))_{n \geq 1}$ est une famille orthogonale
				pour le produit scalaire $\langle f,g\rangle = \int_{-\pi}^{\pi}f(x)g(x)\,dx$.}
			
			\item \question{En déduire que les deux coefficients $b_1$ et $b_2$ peuvent être déterminés
				séparément par
				\[
				b_n = \frac{1}{\pi}\int_{-\pi}^{\pi} f(x)\sin(nx)\,dx.
				\]}
			\reponse{Grâce à l'orthogonalité $\int_{-\pi}^{\pi}\sin(x)\sin(2x)\,dx=0$, les deux
				conditions intégrales se découplent. La première donne
				\[
				\int_{-\pi}^{\pi}f(x)\sin(x)\,dx = b_1\int_{-\pi}^{\pi}\sin^2(x)\,dx = b_1\pi,
				\]
				et la seconde
				\[
				\int_{-\pi}^{\pi}f(x)\sin(2x)\,dx = b_2\int_{-\pi}^{\pi}\sin^2(2x)\,dx = b_2\pi.
				\]
				D'où la formule $b_n = \dfrac{1}{\pi}\displaystyle\int_{-\pi}^{\pi}f(x)\sin(nx)\,dx$
				pour $n=1$ et $n=2$.}
			
			\item \question{Calculer $b_2$.}
			\reponse{On intègre par parties $\int_{-\pi}^{\pi}x\sin(2x)\,dx$, avec $u=x$,
				$v'=\sin(2x)$ :
				\[
				\int_{-\pi}^{\pi}x\sin(2x)\,dx
				= \Bigl[-\frac{x\cos(2x)}{2}\Bigr]_{-\pi}^{\pi}
				+ \frac{1}{2}\int_{-\pi}^{\pi}\cos(2x)\,dx
				= -\pi + 0 = -\pi.
				\]
				Donc $b_2 = \dfrac{1}{\pi}\cdot(-\pi) = -1$.}
			
			\item \question{On admet que, pour tout entier $n\geqslant 1$,
				\[
				\int_{-\pi}^{\pi} x\sin(nx)\,dx = \frac{2\pi(-1)^{n+1}}{n}.
				\]
				En déduire que $b_n = \dfrac{2(-1)^{n+1}}{n}$.}
			\reponse{Directement :
				\[
				b_n = \frac{1}{\pi}\cdot\frac{2\pi(-1)^{n+1}}{n} = \frac{2(-1)^{n+1}}{n}.
				\]
				On vérifie : $b_1=2$, $b_2=-1$, $b_3=\frac{2}{3}$, etc., cohérent avec les calculs
				précédents.}
			
		\end{enumerate}
		
		%-------------------------------------------------------
		\item \texte{\textbf{Vers une écriture en série.}
			
			On obtient ainsi formellement, pour $x\in{]-\pi,\pi[}$ :
			\[
			x = 2\sin(x)-\sin(2x)+\frac{2}{3}\sin(3x)-\frac{1}{2}\sin(4x)+\cdots
			\]
			
			\emph{Note sur la convergence.} Cette égalité repose sur un théorème de convergence
			(théorème de Dirichlet) qui sera démontré ultérieurement. On ne peut pas établir cette
			convergence à l'aide des seuls critères des séries numériques (comparaison, d'Alembert,
			Leibniz), car les termes $\frac{2(-1)^{n+1}}{n}\sin(nx)$ oscillent simultanément en $n$
			et en $x$, ce qui échappe à ce cadre. On admet donc ici que l'égalité est valide sur
			$]-\pi,\pi[$.}
		
		\begin{enumerate}
			
			\item \question{Comparer cette écriture avec les développements en série entière déjà rencontrés :
				quel est le point commun ? Quelle différence importante apparaît ici ?}
			\reponse{\textbf{Point commun :} dans les deux cas, on écrit une fonction comme somme
				infinie de fonctions « élémentaires » (puissances de $x$ pour les séries entières,
				sinus et cosinus ici).
				
				\textbf{Différence :} les séries entières convergent sur un intervalle centré en $0$
				et les fonctions approchées ne sont généralement pas périodiques. Ici, toutes les fonctions
				de la somme sont $2\pi$-périodiques, et la somme approche une fonction périodique.
				De plus, la convergence n'est pas uniforme sur $[-\pi,\pi]$ (elle chute aux bords $\pm\pi$).}
			
			\item \question{En évaluant formellement l'égalité précédente en $x=\dfrac{\pi}{2}$,
				montrer que l'on retrouve la série numérique
				\[
				1-\frac{1}{3}+\frac{1}{5}-\frac{1}{7}+\cdots = \frac{\pi}{4}.
				\]}
			\reponse{En $x=\dfrac{\pi}{2}$ :
				\begin{align*}
					\frac{\pi}{2}
					&= 2\sin\!\Bigl(\frac{\pi}{2}\Bigr)
					-\sin(\pi)
					+\frac{2}{3}\sin\!\Bigl(\frac{3\pi}{2}\Bigr)
					-\frac{1}{2}\sin(2\pi)
					+\frac{2}{5}\sin\!\Bigl(\frac{5\pi}{2}\Bigr)
					-\cdots \\
					&= 2\cdot 1 - 0 + \frac{2}{3}\cdot(-1) - 0 + \frac{2}{5}\cdot 1 - \cdots \\
					&= 2\Bigl(1 - \frac{1}{3} + \frac{1}{5} - \frac{1}{7}+\cdots\Bigr).
				\end{align*}
				En divisant par $2$ :
				\[
				1-\frac{1}{3}+\frac{1}{5}-\frac{1}{7}+\cdots = \frac{\pi}{4}.
				\]
				Il s'agit de la formule de Leibniz-Madhava, ici retrouvée comme cas particulier
				d'une série de Fourier.}
			
		\end{enumerate}
		

		
	\end{enumerate}
}